\documentclass[11pt,a4paper]{article}

% Packages
\usepackage[utf8]{inputenc}
\usepackage[T1]{fontenc}
\usepackage{geometry}
\usepackage{graphicx}
\usepackage{booktabs}
\usepackage{hyperref}
\usepackage{xcolor}
\usepackage{listings}
\usepackage{amsmath}
\usepackage{enumitem}
\usepackage{titlesec}
\usepackage{fancyhdr}

% Page setup
\geometry{margin=1in}
\pagestyle{fancy}
\fancyhf{}
\rhead{Protein-Ligand Binding Prediction}
\lhead{Progress Summary}
\rfoot{Page \thepage}

% Colors
\definecolor{darkblue}{RGB}{0,51,102}
\definecolor{lightgray}{RGB}{240,240,240}

% Hyperref setup
\hypersetup{
    colorlinks=true,
    linkcolor=darkblue,
    urlcolor=darkblue,
}

% Title formatting
\titleformat{\section}{\Large\bfseries\color{darkblue}}{\thesection}{1em}{}
\titleformat{\subsection}{\large\bfseries}{\thesubsection}{1em}{}

\title{\textbf{Protein-Ligand Binding Site Prediction}\\[0.5em]\Large Progress Summary}
\author{Project Collaboration Report}
\date{January 12, 2026}

\begin{document}

\maketitle

\begin{abstract}
This document summarizes the progress made in developing machine learning models for protein-ligand binding site prediction using the PDBbind database. The project explored multiple approaches, including Random Forest, Fully Connected Neural Networks (FCNN), Transformer classifiers, and Context-Aware Transformers, each with varying levels of success in predicting amino acid-level binding sites.
\end{abstract}

\tableofcontents
\newpage

%------------------------------------------------------------------------------
\section{Project Overview}

\subsection{Objective}
The main goal of this project is to develop computational methods for predicting protein-ligand binding sites at the residue (amino acid) level. Given a protein sequence and ligand properties, the models predict which amino acids are likely to be involved in binding interactions.

\subsection{Data Source}
The project uses the \textbf{PDBbind database}, a comprehensive collection of protein-ligand complexes with experimentally determined structures and binding affinities.

\subsection{Problem Formulation}
\begin{itemize}
    \item \textbf{Input}: Amino acid features + global protein features + ligand properties
    \item \textbf{Output}: Binary classification (binding / non-binding) per residue
    \item \textbf{Challenge}: Highly imbalanced data (binding sites are rare)
\end{itemize}

%------------------------------------------------------------------------------
\section{Feature Engineering}

\subsection{Baseline Feature Set (46 Dimensions)}
The initial feature extraction pipeline generates 46-dimensional feature vectors:

\begin{table}[h]
\centering
\begin{tabular}{@{}llll@{}}
\toprule
\textbf{Index Range} & \textbf{Feature Group} & \textbf{Description} & \textbf{Dims} \\ \midrule
0--19 & Residue Identity & One-Hot Encoding (20 amino acids) & 20 \\
20--21 & Residue PhysChem & Charge and Hydrophobicity & 2 \\
22--41 & Global Protein & Amino Acid Composition (AAC) & 20 \\
42--45 & Ligand Features & LogP, MW, H-Donors, H-Acceptors & 4 \\ \bottomrule
\end{tabular}
\caption{Baseline 46-dimensional feature vector composition}
\end{table}

\subsection{Enhanced Feature Set (466 Dimensions)}
The enhanced pipeline adds sophisticated protein sequence descriptors:

\begin{table}[h]
\centering
\begin{tabular}{@{}lllll@{}}
\toprule
\textbf{Index Range} & \textbf{Feature Group} & \textbf{Description} & \textbf{Dims} & \textbf{Source} \\ \midrule
0--19 & Residue Identity & One-Hot Encoding & 20 & Custom \\
20--21 & Residue PhysChem & Charge, Hydrophobicity & 2 & Custom \\
22--41 & AAC & Amino Acid Composition & 20 & pypropel \\
42--441 & DAC & Dipeptide Composition & 400 & pypropel \\
442--461 & AVEANF & Amphiphilic Pseudo-AA & 20 & pypropel \\
462--465 & Ligand Features & LogP, MW, H-Donors, H-Acceptors & 4 & RDKit \\ \bottomrule
\end{tabular}
\caption{Enhanced 466-dimensional feature vector composition}
\end{table}

%------------------------------------------------------------------------------
\section{Model Development}

\subsection{Random Forest Classifier}

\subsubsection{Approach}
The first model implemented was a Random Forest classifier, chosen for its:
\begin{itemize}
    \item Robustness to feature scaling
    \item Built-in feature importance
    \item Ability to handle imbalanced data
\end{itemize}

\subsubsection{Results}

\begin{table}[h]
\centering
\begin{tabular}{@{}lllllll@{}}
\toprule
\textbf{Configuration} & \textbf{Data Size} & \textbf{ROC-AUC} & \textbf{Accuracy} & \textbf{Recall} & \textbf{Precision} \\ \midrule
Baseline (46 dims) & 5 Batches (~200k) & \textbf{0.8013} & 0.8330 & 0.54 & 0.20 \\
Enhanced (466 dims) & 5 Batches (~227k) & 0.7248 & \textbf{0.8687} & 0.36 & 0.15 \\
Enhanced (466 dims) & Full (~569k) & \textbf{0.7961} & 0.8387 & \textbf{0.55} & 0.19 \\ \bottomrule
\end{tabular}
\caption{Random Forest model comparison (December 27, 2025)}
\end{table}

\subsubsection{Key Findings}
\begin{enumerate}
    \item \textbf{Baseline Effectiveness}: The simple 46-dim feature set is robust (AUC $\sim$0.80) even with limited data.
    \item \textbf{Curse of Dimensionality}: Enhanced features underperformed when data was limited (AUC 0.72 vs 0.80).
    \item \textbf{Data Scale Matters}: Full dataset allowed enhanced features to recover performance.
    \item \textbf{Performance Ceiling}: More features did not strictly outperform baseline, suggesting need for more complex models.
\end{enumerate}

%------------------------------------------------------------------------------
\subsection{Fully Connected Neural Network (FCNN)}

\subsubsection{Approach}
A simple multi-layer perceptron was implemented to establish a neural network baseline, using the enhanced 466-dimensional features.

\subsubsection{Training Configuration}
\begin{itemize}
    \item Focal Loss to handle class imbalance
    \item Adam optimizer with learning rate scheduling
    \item Early stopping based on validation AUC
\end{itemize}

%------------------------------------------------------------------------------
\subsection{Transformer Classifier}

\subsubsection{Architecture}
The Transformer-based classifier treats the 466-dimensional feature vector as a sequence of feature patches:

\begin{enumerate}
    \item \textbf{Input Projection}: Features divided into patches ($\sim$20 features each)
    \item \textbf{Positional Encoding}: Sinusoidal encoding for patch positions
    \item \textbf{CLS Token}: Learnable classification token prepended
    \item \textbf{Transformer Encoder}: Self-attention layers learn feature interactions
    \item \textbf{Classification Head}: MLP with sigmoid output
\end{enumerate}

\subsubsection{Variants}
\begin{itemize}
    \item \textbf{BindingTransformer}: Full architecture with patch-based input
    \item \textbf{BindingTransformerLite}: Lighter version with direct projection
\end{itemize}

%------------------------------------------------------------------------------
\subsection{Context-Aware Transformer}

\subsubsection{Motivation}
To leverage sequence context for binding prediction, treating each prediction as dependent on neighboring amino acids.

\subsubsection{Architecture}
\begin{enumerate}
    \item \textbf{Input}: Window of $2k+1$ amino acids (context size $k$)
    \item \textbf{Learnable Positional Encoding}: For short sequences
    \item \textbf{Self-Attention}: Learns inter-residue relationships
    \item \textbf{Center Token Extraction}: Prediction for center residue
    \item \textbf{Classification Head}: Binary binding probability
\end{enumerate}

\subsubsection{Key Components}
\begin{verbatim}
class ContextTransformer(nn.Module):
    # context_size: Residues on each side of center
    # d_model: Transformer dimension (256)
    # nhead: Attention heads (8)
    # num_layers: Encoder layers (4)
\end{verbatim}

\subsubsection{Alternative: CLS Token Architecture}
Also implemented \texttt{ContextTransformerWithCLS} that uses a learnable CLS token for classification instead of the center token.

%------------------------------------------------------------------------------
\section{Regression Experiments}

\subsection{Log-Transform Distance Regression}
An alternative approach was explored: predicting the distance to the nearest ligand atom as a regression task.

\subsubsection{Approach}
\begin{itemize}
    \item Target: Log-transformed distance to nearest ligand atom
    \item Model: Transformer-based regression architecture
    \item Loss: Mean Squared Error (MSE)
\end{itemize}

\subsubsection{Threshold-Based Classification}
Post-hoc classification by applying distance thresholds to regression predictions.

\subsubsection{Findings}
\begin{itemize}
    \item Model struggled to predict very short distances (actual binding sites)
    \item MSE loss dominated by large distances
    \item Direct classification remains more effective for this task
\end{itemize}

%------------------------------------------------------------------------------
\section{Training Infrastructure}

\subsection{Loss Functions}
\begin{itemize}
    \item \textbf{Focal Loss}: Addresses class imbalance by down-weighting easy examples
    \item \textbf{Weighted BCE}: Class-weighted binary cross-entropy
    \item \textbf{MSE}: For regression experiments
\end{itemize}

\subsection{Data Pipeline}
\begin{itemize}
    \item Batch processing of PDB structures
    \item Protein-level train/val/test split (prevents leakage)
    \item Context window generation for sequence models
    \item On-the-fly data loading for memory efficiency
\end{itemize}

\subsection{Hardware}
\begin{itemize}
    \item Training on Apple Silicon (MPS) when available
    \item CUDA support for GPU clusters
    \item CPU fallback
\end{itemize}

%------------------------------------------------------------------------------
\section{Saved Models}

The following trained models are available:

\begin{table}[h]
\centering
\begin{tabular}{@{}ll@{}}
\toprule
\textbf{Model File} & \textbf{Description} \\ \midrule
\texttt{rf\_model.joblib} & Random Forest classifier \\
\texttt{rf\_model\_enhanced.joblib} & RF with enhanced features \\
\texttt{rf\_model\_enhanced\_full.joblib} & RF on full dataset \\
\texttt{best\_fcnn\_focal.pth} & FCNN with Focal Loss \\
\texttt{best\_transformer\_focal.pth} & Transformer classifier \\
\texttt{best\_transformer\_lite\_focal.pth} & Lite Transformer \\
\texttt{best\_context\_transformer\_focal.pth} & Context-Aware Transformer \\
\texttt{best\_regression\_transformer\_mse.pth} & Distance regression model \\
\texttt{best\_regression\_transformer\_mse\_log.pth} & Log-distance regression \\ \bottomrule
\end{tabular}
\caption{Saved model checkpoints}
\end{table}

%------------------------------------------------------------------------------
\section{Code Organization}

\subsection{Source Code Structure}
\begin{verbatim}
src/
|-- dataset.py              # Base dataset loader
|-- dataset_context.py      # Context window dataset
|-- dataset_regression.py   # Regression target dataset
|-- features.py             # Feature extraction
|-- loss.py                 # Custom loss functions (Focal)
|-- model.py                # Base model definitions
|-- model_context_transformer.py   # Context-aware models
|-- model_regression.py     # Regression models
|-- model_rf.py             # Random Forest wrapper
|-- model_transformer.py    # Transformer classifiers
|-- pipeline.py             # Training pipeline
|-- pipeline_regression.py  # Regression training
+-- utils.py                # Utility functions
\end{verbatim}

\subsection{Training Scripts}
\begin{itemize}
    \item \texttt{train\_fcnn.py} -- FCNN training
    \item \texttt{train\_transformer.py} -- Transformer training
    \item \texttt{train\_context\_transformer.py} -- Context-aware training
    \item \texttt{train\_regression.py} -- Distance regression training
\end{itemize}

\subsection{Evaluation Scripts}
\begin{itemize}
    \item \texttt{evaluate\_fcnn.py}
    \item \texttt{evaluate\_transformer.py}
    \item \texttt{evaluate\_context\_transformer.py}
    \item \texttt{evaluate\_regression.py}
\end{itemize}

%------------------------------------------------------------------------------
\section{Conclusions and Future Directions}

\subsection{What Worked}
\begin{enumerate}
    \item \textbf{Baseline features} (46 dims) provide strong performance with simple models
    \item \textbf{Focal Loss} effectively handles class imbalance
    \item \textbf{Protein-level data splitting} prevents data leakage
    \item \textbf{Context-aware modeling} captures sequence dependencies
\end{enumerate}

\subsection{Challenges}
\begin{enumerate}
    \item \textbf{High-dimensional features} require more data to avoid overfitting
    \item \textbf{Class imbalance} remains a fundamental challenge
    \item \textbf{Precision-recall trade-off} difficult to optimize
\end{enumerate}

\subsection{Recommended Next Steps}
\begin{enumerate}
    \item \textbf{Feature Selection}: Apply feature selection on the 466-dim set to remove noise
    \item \textbf{Gradient Boosting}: Explore XGBoost/LightGBM for high-dimensional features
    \item \textbf{Structure Features}: Incorporate 3D structural information
    \item \textbf{Pre-trained Embeddings}: Use protein language models (ESM, ProtTrans)
    \item \textbf{Multi-task Learning}: Joint prediction of binding sites and affinity
\end{enumerate}

%------------------------------------------------------------------------------
\section{Timeline}

\begin{tabular}{@{}ll@{}}
\toprule
\textbf{Date} & \textbf{Milestone} \\ \midrule
Dec 27, 2025 & Feature extraction pipeline completed \\
Dec 27, 2025 & Random Forest PoC with baseline \& enhanced features \\
Dec 27, 2025 & FCNN implementation and training \\
Dec 27, 2025 & Transformer classifier development \\
Jan 5--6, 2026 & Log-transform regression experiments \\
Jan 12, 2026 & Context-Aware Transformer implementation and evaluation \\ \bottomrule
\end{tabular}

\vspace{2em}
\hrule
\vspace{1em}
\textit{Document generated on January 12, 2026}

\end{document}
